\section{Scope, results, and relation to the repository}\label{sec:scope}

In this article, \emph{series reversal} means \emph{compositional reversion}:
given a relation \(y=f(t)\), determine and justify a local inverse \(t=g(y)\).
It does not mean reversing the order of coefficients or taking \(1/f\).
The central issue is that an inverse may need a different asymptotic scale
from the original function.  The elementary example
\[
 y=\ee^{-S/t},\qquad
 t=-\frac{S}{\Log y+2\pi\ii j},\qquad j\in\Z,
\]
already exchanges an exponential scale for a logarithmic covering.
If the exponential is multiplied by a power of \(t\), Lambert \(W\)
enters.  If a leading correction is \(Q^m\), a Puiseux root enters before
the logarithm.  If the original coefficients are divergent, a choice of
Borel summation direction enters as well.

These mechanisms are often discussed separately.  Here they are joined in
an exact class arising from modular \(q\)-products.  The resulting theorems
include branch data and error bounds; they do not depend on promoting an
uncontrolled formal substitution to an analytic identity.

\subsection{What is proved}

The main results have five distinct roles.
\begin{enumerate}[leftmargin=*,itemsep=4pt]
\item \Cref{thm:cyclotomic-borel,thm:direction-counterexample} give an
explicit Borel transform for the fixed-argument cyclotomic \(q\)-Pochhammer
logarithm and a counterexample to one directional-equality assertion in
the repository.  The general resurgence mechanism is classical; the
point here is to state and prove the correctly normalized specialization.
\item \Cref{thm:core-inverse,thm:finite-endpoint} give two practical
reversion engines.  One treats an exact dominant core with a convergent
exponential tail.  The other treats a regular finite endpoint perturbed by
\(\ee^{-S/t}\), and identifies the sharp scale and all leading coefficient
poles.
\item \Cref{thm:cusp-normal-form,thm:flat-inverse} give the exact
normal form at every rational cusp and a convergent inverse when all
perturbative core terms cancel.
\item \Cref{thm:three-factors,thm:four-factors} construct arbitrary
finite ramification with minimal numbers of factors: three eta factors
for the two eta balances, four Euler-product factors for all three raw
product balances.
\item \Cref{thm:qgamma-completion} and the subsequent half-argument
calculation apply the theory to rational special values of \(q\)-gamma,
both before and after removal of their analytic background.
\end{enumerate}

A precise negative conclusion is part of the result:
a Borel ray containing no pole need not produce the same sum as another
such ray.  A precise positive conclusion is equally important:
after fixing the summation and logarithm sheets, an inverse can be
constructed with every exponential contribution retained and with a
geometric remainder.

\subsection{Existing results that must be credited}

The source comparison is pinned to commit
\texttt{6cb1d86f19ebf9f54c8d985b141b84233b6eab66}.
The \repo{} transseries volume and companion already contain exact-core
inversion, Lagrange--B\"urmann formulas, exponential corrections, and
positive-real \(q\)-product completions \cite{repo-trans,repo-companion}.
The support-controlled reversion and moving-fold reports already discuss
finite action grids, grouped resonances, and analytic critical charts
\cite{repo-reversion,repo-fold}.  The present argument does not claim
those principles as new.

The relevant source base also extends beyond the Transseries directory.
The \(q\)-Pochhammer monograph contains complex fixed-argument expansions,
cyclotomic asymptotic base inverses, exact eta completions at \(q=1,-1\),
ordinary finite ramification, and regular \(q\)-gamma base inversion
\cite{repo-qmono}.  It would therefore be incorrect to describe the
repository as confined to positive real \(q\).  Its frontier chapter
contains the named conjecture
\texttt{conj:cyclotomic-resurgent-inverse}, whose directional-equality
clause we test below.

The literature also goes substantially beyond a purely formal calculus.
Van der Hoeven proves formal compositional inversion in an appropriate
transseries field and develops a separate complex theory with
composition-domain restrictions \cite{vdh-book,vdh-complex}.
Edgar studies composition, recursion, and convergence
\cite{edgar}.  Marmi and Sauzin treat Borel transforms at roots-of-unity
resonances for a general cohomological equation
\cite[Sections 4.2--4.3]{marmi-sauzin}.  Their framework includes the
Lambert series used here.  Explicit quantum-dilogarithm Borel poles and
Stokes jumps are developed by Garoufalidis and Kashaev
\cite{garoufalidis-kashaev}; more recent \(q\)-Pochhammer modular
resurgence is treated by Fantini and Rella \cite{fantini-rella}.
Sauzin establishes nonlinear resurgent substitution and inversion under
the required analytic estimates \cite{sauzin}.

Our contribution is an explicit development within this established
framework: a correction of a specific repository assertion, exact
logarithmic cusp inversion with arithmetic coefficient generation,
sharp factor-count constructions, and the finite-endpoint feedback
analysis for modular \(q\)-gamma corrections.  These are proved below.
A worldwide priority claim, or the resolution of unrestricted complex
transseries inversion, would go beyond the evidence.

\subsection{Notation and domain conventions}

We use \(\N=\{0,1,2,\ldots\}\).  Put
\[
 P(z)=\prod_{n=1}^{\infty}(1-z^n),\qquad |z|<1.
\]
This is holomorphic and nonzero in the unit disk, and its logarithm is
the unique analytic logarithm with \(\log P(0)=0\).
The eta function is
\[
 \eta(\tau)=\ee^{\pi\ii\tau/12}P(\ee^{2\pi\ii\tau}),
 \qquad \Im\tau>0.
\]
For an integer \(n\ge1\), let
\(\sigma_{-1}(n)=\sum_{d\mid n}d^{-1}\).
A root of unity is written
\(\zeta=\ee^{2\pi\ii h/k}\), with \(k\ge1\) and \((h,k)=1\);
the cusp \(q=1\) is represented by \(h=0,k=1\).
We use
\[
 q=\zeta\ee^{-t},\qquad \Re t>0,\qquad
 \sectors=\{t:0<|t|<t_*,\ |\arg t|\le\pi/2-\delta\}.
\]
Here \(0<\delta<\pi/2\), and every asymptotic statement on
\(\sectors\) is uniform on the displayed closed subsector unless
explicitly qualified.

The symbol \(\Log\) denotes a specified branch or logarithm lift.
A logarithm of a holomorphic unit equal to \(1\) at the origin always
means its germ with value \(0\) there.  Thus a logarithm sheet is not
silently reset during continuation.  For complex \(q=\ee^{-t}\) and
rational \(x\), the definition of \(q^x\) is \(\ee^{-tx}\);
this is the branch used in the \(q\)-gamma applications.

\section{Complex reversion: formal data and analytic realization}
\label{sec:complex-framework}

\subsection{Three different meanings of a transseries inverse}

There are three levels of assertion.
A \emph{formal inverse} is an identity in a specified complete algebra
of series, where the relevant substitutions are defined.
An \emph{asymptotic inverse} is a sectorial function with controlled
remainders after truncation.
An \emph{exact convergent inverse representation} is an identity
obtained by substituting small functions into a convergent analytic
series.  Such a representation is stronger than its asymptotic expansion,
but does not supply all possible divergent or oscillatory transseries.

For example, if \(U(z_1,z_2)\) is analytic near \((0,0)\), then
\[
 U\bigl(\ee^{-\alpha/t},\ee^{-\beta/t}\bigr)
 =\sum_{p,q\ge0}u_{p,q}\ee^{-(p\alpha+q\beta)/t}
\]
converges whenever both arguments lie in a suitable polydisk.
It is an asymptotic transseries on a sector only when their real
exponential rates decay there.  If \(p\alpha+q\beta\) agrees for two
different pairs, their coefficients must be added before the leading
term is identified.  No denominator involving an action difference is
needed for this addition.

\begin{proposition}[Finite complex actions]\label{prop:finite-actions}
Let \(F(u,\mathbf z)\) be holomorphic near \((0,\mathbf0)\), with
\(F(0,\mathbf0)=0\) and \(F_u(0,\mathbf0)\ne0\).  There is a unique
holomorphic germ \(u=U(\mathbf z)\) with \(U(\mathbf0)=0\) and
\(F(U(\mathbf z),\mathbf z)=0\).
If \(z_j(t)=\ee^{-\alpha_j/t}\), and a closed sector satisfies
\[
 \Re(\alpha_j/t)\ge c/|t| \quad (1\le j\le r)
\]
for a fixed \(c>0\), then \(U(\mathbf z(t))\) is an exact convergent
sectorial transseries.  Its action support is contained in
\(\N\alpha_1+\cdots+\N\alpha_r\), and equal actions are grouped by
finite sums at each bounded exponential height.
\end{proposition}
\begin{proof}
The holomorphic implicit-function theorem gives \(U\) and a polydisk
of convergence.  Shrinking the sector radius puts all \(z_j(t)\)
in any smaller such polydisk.  Absolute convergence permits grouping.
On a fixed interior ray, the real part of
\(\sum n_j\alpha_j/t\), divided by \(1/|t|\), is at least
\(c\sum n_j\).  Consequently only finitely many multiindices lie below
a fixed height.  This proves both the finite grouping assertion and
the usual asymptotic interpretation of degree truncations.
\end{proof}

The finite-action result is standard analytic infrastructure.  It remains
useful because it says exactly which operation is legitimate.  An
arbitrary countable action family would additionally require control of
accumulation, coefficient summability, and substitution.  An arbitrary
complex direction may not meet a common decay cone at all.

\subsection{Why a global complex claim fails}

The size of \(\ee^{-1/t}\) changes when
\(\Re(1/t)\) changes sign.  The two functions
\[
 f_0(t)=t,\qquad f_1(t)=t+\ee^{-1/t}
\]
have the same algebraic asymptotic series on every closed right
subsector, but different inverses there.  Thus an algebraic expansion
does not determine exponentially small inverse information.
The exact normalizations in the following sections are what supply it.

A finite-point critical example is
\(f(t)=t^m\).  Its inverse requires an \(m\)-th root; it is not an
ordinary Taylor series at the critical value when \(m>1\).
At a flat cusp, \(\ee^{-1/t}\), even a Puiseux scale alone is
insufficient: a logarithm is unavoidable.
For arbitrary complex exponential actions, a sector may also contain
oscillatory pairs of comparable size.  A total real ordering of
monomials cannot simply be imported unchanged into that situation.
These distinctions are consistent with the composition restrictions
in the complex transseries literature \cite{vdh-complex}.

\subsection{A formal coefficient ring is not a summation theorem}

For coefficients in a complete formal algebra, an equation
\[
 u=-\Phi(u,\mathbf z),\qquad \Phi(0,\mathbf0)=0,
\]
with a suitably invertible linearization can be solved by coefficient
recursion.  In a \(\mathbf z\)-adic algebra, every coefficient uses
finitely many preceding coefficients.  That fact proves formal
existence and uniqueness.  It does not give bounds after substituting
a physical exponential scale.

For a resurgent extension, it is also insufficient that each
coefficient is individually resurgent.  An analytic series in a
second variable must have geometric coefficient bounds in the
relevant resurgent seminorms.  Sauzin's nonlinear implicit theorem
provides the appropriate setting \cite{sauzin}.  Our convergent
cusp theorems meet their analytic needs directly through ordinary
holomorphic estimates; our Borel calculations specify the
summation data rather than assuming it.

